%% depth-completion-supplement.tex
%% Separately authored completion of the historical Sigma-one proof outline.

% Necessity: the body uses both external and formalized Sigma-one completeness.
% Definition: this file supplies the bounded-formula induction, normal-form
% conversion, proof-code constructors, substitution invariant, and internal PA
% verification omitted by the historical outline.
% Standing: it is a reciprocal proof annex for one body theorem.
% Nonredundancy: exact inherited proofs stay inline in the body; this item proves
% only the separately authored Sigma-one-completeness theorem.
% Use: annex/p05-route inputs this file once.
\noindent\textit{Reader note.}  This separately authored annex proof completes
Theorem~\ref{thm:p05-sigma-one-completeness}. Its internal clause uses the
literal selected proof predicate of
Definition~\ref{def:p05-selected-direct-pa-graph} and the specified builders
of \cref{thm:p05-fixed-bounded-builders,thm:p05-original-sigma-evidence}.

% BEGIN HISTORICAL SLICE 16-34
\begin{annexitem}{annex:p5.sigma.1}{thm:p05-sigma-one-completeness}{External and formalized Sigma-one completeness}
Fix a closed \(\Sigma^0_1\) sentence \(\sigma\) in its original syntax,
using the closure-style definition in
Table~\ref{tab:p04-arithmetical-hierarchy}. An existential witness below a
bounded universal may depend on the bounded variable. We first prove
bounded-instance completeness and then encode all finitely many dependent
witnesses in one evidence number. The fixed marked negation-normal-form
conversion determines its evidence; the concluding proof has the original
sentence as its root.

For a term \(t(\vec x)\) and a tuple \(\vec n\), let
\(\operatorname{val}_t(\vec n)\) be its value in \(\mathcal N\).
Structural recursion on \(t\) constructs a \(\mathsf Q\)-proof of
\[
 t(\overline{\vec n})=
 \overline{\operatorname{val}_t(\vec n)}.                            \tag{2}
\]
For variables and \(0\) this is immediate.  The successor case uses equality
congruence.  The addition and multiplication cases append, respectively, the
axioms \(x+0=x\), \(x+Sy=S(x+y)\), \(x\times0=0\), and
\(x\times Sy=(x\times y)+x\) once for each digit of the evaluated second
argument.  Thus the map from \(\vec n\) to the code of (2) is primitive
recursive, not merely computable by an unbounded proof search.

We now prove, simultaneously for every bounded formula
\(\theta(\vec x)\), the two numeral-instance clauses
\[
\begin{array}{ll}
 \mathcal N\entails\theta(\vec n)
   &\Longrightarrow
   \mathsf Q\derives\theta(\overline{\vec n}),\\[2pt]
 \mathcal N\not\entails\theta(\vec n)
   &\Longrightarrow
   \mathsf Q\derives\neg\theta(\overline{\vec n}).
\end{array}                                                           \tag{3}
\]
The induction also constructs the proof codes primitive recursively from
\(\vec n\).

For an atomic equality, evaluate its two terms by (2).  Equal numeral values
give the equality by transitivity; unequal values give its negation by repeated
successor injectivity and \(Sx\ne0\). Bottom is refuted by the fixed
logical derivation of \(\neg\bot\). For comparisons retain the two actual
abbreviations
\[
 u\leq v\defeq\exists z\,(z+u=v),\qquad
 u<v\defeq\exists z\,(u+Sz=v).
\]
True numeral comparisons use their standard difference witness and (2).
False weak comparisons follow from \cref{lem:q-bounded-cases} and
the refutations of the resulting numeral equalities.

Here is the strict enumeration needed both for false numeral comparisons
and strict bounded quantifiers. Metainduction on \(m\) constructs a
\(\mathsf Q\)-proof that \(y<\bar m\) implies the disjunction of
\(y=\bar j\) for \(j<m\), with empty disjunction \(\bot\).
At zero, \(y+Sz=0\) contradicts \(y+Sz=S(y+z)\) and \(Sx\ne0\).
At \(m+1\), \(y+Sz=S\bar m\) gives \(y+z=\bar m\).
Split \(z\) using the zero-or-successor axiom. If \(z=0\), then
\(y=\bar m\). If \(z=Sw\), then \(y+Sw=\bar m\), and the
previous strict enumeration gives one of \(y=\bar j\) for \(j<m\).
These are exactly all cases below \(m+1\). Each step uses a fixed finite
proof template, so this construction is primitive recursive in \(m\).
For a false strict numeral comparison, substitute its first numeral in
this enumeration and refute every resulting equality. Finally the Boolean
cases append the corresponding fixed propositional derivations to the
recursively constructed proofs.

For a bounded existential
\(\exists y\leq t(\vec x)\,\eta(y,\vec x)\), evaluate the bound to \(m\).
If the formula is true, bounded search returns the least \(j\leq m\) for which
\(\eta(j,\vec n)\) is true; the positive induction clause, the numeral proof
of \(\bar j\leq\bar m\), and existential introduction give the required
proof.  If it is false, the negative induction clause supplies proofs of
\(\neg\eta(\bar j,\overline{\vec n})\) for every \(j\leq m\).
Lemma~\ref{lem:q-bounded-cases} replaces an arbitrary
\(y\leq\bar m\) by the finite disjunction
\(y=\bar0\lor\cdots\lor y=\bar m\); reasoning by cases proves the negation
of the bounded existential.  The bounded universal case is dual: use all
positive instance proofs and the same finite case split when it is true, and
use the least counterexample when it is false. For a strict bound use the
cases \(j<m\); when \(m=0\), the strict guard is contradictory, so its
universal is proved and its existential refuted without a matrix instance.
A weak zero bound has the single case \(j=0\). Each inner bound is
evaluated after its outer parameters and witnesses have been substituted.
All loops have the resulting finite bounds, so this completes both the
structural induction and the primitive-recursive construction in (3).

For the internal bounded-instance assertion, take the actual descriptions
\(d_\theta^+,d_\theta^-\) of
\cref{def:p05-fixed-bounded-builders}. These pack the specified raw trees
\(R_\theta^+,R_\theta^-\), with their respective success guards and zero
failure values. List the free indices of \(\theta\) as
\(i_1<\cdots<i_k\). Let \(s_\theta^+\) be the finite composition that
starts with \(\ulcorner\theta\urcorner\) and, in that order, applies
the selected \(\mathsf{sub}\) description at index \(i_j\) with the
numeral description \(\nu(x_j)\). Put
\(s_\theta^-=\mathsf{Neg}\circ s_\theta^+\), using the fixed
\(N_7\) constructor. These are exactly the descriptions denoted by
\(\theta^\nu(\vec x)\) and
\(\mathsf{Neg}(\theta^\nu(\vec x))\) in the bounded-builder theorem.
For \(k=0\), the first description is the constant original code.
The graph theorem \cref{thm:p05-pa-pr-graphs} gives existence and uniqueness
for each of these separately fixed descriptions, inside \(\PA\).

Theorem~\ref{thm:p05-fixed-bounded-builders}, proved in
\cref{annex:p5.fixed-bounded-builders}, therefore gives the following
formulas, with universal closure in \(\vec x\):
\[
\begin{aligned}
 \PA&\derives\theta(\vec x)\to\exists p\exists q\,\\
 &\quad(\Gamma_{d_\theta^+}(\vec x,p)\land
   \Gamma_{s_\theta^+}(\vec x,q)\land
   \mathsf{Prf}^{19}_{\PA}(p,q)),\\
 \PA&\derives\neg\theta(\vec x)\to\exists p\exists q\,\\
 &\quad(\Gamma_{d_\theta^-}(\vec x,p)\land
   \Gamma_{s_\theta^-}(\vec x,q)\land
   \mathsf{Prf}^{19}_{\PA}(p,q)).
\end{aligned}
\tag{4}
\]
Here \(\mathsf{Prf}^{19}_{\PA}(p,q)\) is literally
\(\Gamma_{d_{19}}(p,q,\bar1)\). For example, uniqueness of the two
positive graph outputs turns the first line into the corresponding
implication for any supplied \(p,q\) satisfying those graphs. The cited
proof establishes the Boolean truth test, the safe-tree invariant, the
literal root after substitution, and all sentence and open-axiom factors
of that checker after packing. Its bounded folds retain dependent
parameters and both strict and weak zero-bound cases. Thus (4) uses the
internal equations and verification of these particular descriptions.

We now pass from bounded formulas to the full closure-style
\(\Sigma^0_1\) class. For every fixed \(\Sigma^0_1\)-formula
\(\varphi(\vec x)\), take precisely the fixed marked normal form
\(\rho=\mathcal N^+\varphi\) in
\cref{def:p05-signed-evidence-builders}. Define
\(E_\varphi(\vec x,e)\) to be an alias for the existing
\(\mathsf{Ev}_\varphi(\vec x,e)\) of
\cref{thm:p05-pa-sigma-evidence}: it is
\(\mathcal E_\rho(\vec x,e)\), with the very same fresh evidence
variable, auxiliary bound variables, conjunction associations and original
comparison binders as in
\cref{def:p05-positive-evidence-formation}. This is an assignment of a
fixed bounded formula, not an additional object-language predicate family.
For clarity, its clauses on the positive formation \(\rho\) are as
follows; the bounded-leaf clause always has priority.
\begin{enumerate}
  \item If the current \(\rho\) is bounded, use
        \(E_\rho(\vec x,e)\defeq(e=0\land\rho(\vec x))\).
        In particular a bounded original \(\varphi\) has the leaf
        \(e=0\land\mathcal N^+\varphi\).
  \item For \(\rho=\psi\land\chi\), use
        \[
        \exists u\leq e\,\exists v\leq e\,
          (\mathsf P(u,v,e)\land E_\psi(\vec x,u)
                             \land E_\chi(\vec x,v)).
        \]
        For \(\psi\lor\chi\), use
        \[
        \exists d<\bar2\,\exists v\leq e\,
          [\mathsf P(d,v,e)\land
            ((d=0\land E_\psi(\vec x,v))\lor
             (d=\bar1\land E_\chi(\vec x,v)))].
        \]
  \item For \(\rho=\exists y\,\psi(y,\vec x)\), use
        \[
        \exists y\leq e\,\exists v\leq e\,
          (\mathsf P(y,v,e)\land E_\psi(y,\vec x,v)).
        \]
        For a marked existential replace \(y\leq e\) by its original
        guard \(y\mathrel\diamond t(\vec x)\), where
        \(\diamond\in\{<,\leq\}\); retain \(v\leq e\).
  \item For
        \(\rho=\forall y\mathrel\diamond t(\vec x)\,\psi(y,\vec x)\), use
        \[
        E_\rho(\vec x,e)\defeq
        \forall y\mathrel\diamond t(\vec x)\,
        \exists v\leq e\,
          (\mathsf B(e,y,v)\land E_\psi(y,\vec x,v)),              \tag{5}
        \]
        where \(\mathsf B\) is the original bounded beta-entry formula of
        Lemma~\ref{lem:goedel-beta}.
\end{enumerate}
The weak guards \(u\leq e,v\leq e\) are the actual guards in
\(\mathsf{Ev}\); they are not replaced by \(u<Se,v<Se\).
Pairing components and beta entries are at most their respective code,
including at code zero. Every newly written quantifier is marked bounded.
Consequently each \(E_\varphi\) is bounded. The evidence number at a
bounded universal stores a possibly different existential-evidence object
for each bounded index; it is used in (5) and the two implications below.

Two metainductions on the fixed positive formation \(\rho\), followed
by its fixed logical equivalence to \(\varphi\), give
\[
 \mathsf Q\derives
   E_\varphi(\vec x,e)\to\varphi(\vec x),                          \tag{6}
\]
and
\[
 \PA\derives
   \varphi(\vec x)\to\exists e\,E_\varphi(\vec x,e).                \tag{7}
\]
Both displays have universal closure in their free numerical variables.
For (6), a bounded leaf follows by conjunction elimination. A conjunction
supplies evidence for both children; a disjunction supplies one of its two
tagged alternatives. Apply the child implications and the corresponding
propositional introductions. An existential supplies a witness and its
matrix evidence; apply the child implication and introduce that same
witness, retaining its original guard in a marked case. In the universal
case, (5) supplies matrix evidence for each \(y\) in its original range;
the child implication, guard discharge and universal introduction give
that marked universal. These steps use only first-order logic and so are
available in \(\mathsf Q\). They yield \(E_\varphi\to\rho\).
Each fixed negation-normal-form rewrite, including implication elimination,
double negation, De Morgan laws and quantifier duality, is a classical
first-order equivalence. Expanding a marked quantifier gives its bounded
duality while retaining its bound and comparison binder; no rewrite moves
a formula across a binder. Consequently \(\mathsf Q\derives
\rho\leftrightarrow\varphi\), and composition proves (6) for the
original syntax. Auxiliary variables are fresh, and an otherwise unused
assumption can be used by conjunction and projection before discharge.

The nonlogical existence step at a bounded universal in (7) is finite
evidence collection. For every formula \(R(i,u,\vec a)\), full induction
in \(\PA\) proves the scheme
\[
 \forall i<n\,\exists u\,R(i,u,\vec a)
 \ \to\
 \exists w\,\forall i<n\,\exists u\leq w\,
   (\mathsf B(w,i,u)\land R(i,u,\vec a)).                         \tag{8}
\]
This is \cref{lem:p05-pa-finite-collection}. Its proof inducts on \(n\)
with all parameters retained. At zero take \(w=0\). At \(Sn\), use
the induction hypothesis and choose the new witness at \(n\); the
internally proved extension lemma \cref{lem:p05-pa-beta-extension}
appends it while preserving the first \(n\) entries. The internal
pairing and beta lemma \cref{lem:p05-pa-pair-beta} bounds every entry by
the new code. These two cases prove (8) inside \(\PA\). Substitution
of \(Sm\) for \(n\) gives its weak-bound variant, including the
single witness at a weak zero bound.

For (7), a bounded leaf uses evidence zero. Pair the two child evidence
values at a conjunction, or the appropriate tag and child evidence at a
disjunction. At an existential pair its witness with the evidence of its
matrix; pairing bounds supply the weak evidence guards, and a marked
witness retains its original comparison. At a bounded universal use (8)
after the child existence implication, with
\(R(y,v,\vec x)=E_\psi(y,\vec x,v)\). Choose
\(n=t(\vec x)\) for a strict bound and \(n=St(\vec x)\) for a
weak bound. A strict zero bound uses the empty code and needs no matrix
witness. Inner witnesses and bounds keep all their outer parameters, so
this collection covers dependent witnesses. These are precisely the
evidence-existence derivations of
\cref{thm:p05-pa-sigma-evidence,thm:p05-evidence-projections} and
\cref{annex:p5.evidence-projections}. Composing with
\(\varphi\leftrightarrow\rho\) gives (7).

For the internal proof record, take the specified signed description
\(\widehat d_\varphi^+\) of
\cref{def:p05-signed-evidence-builders}. Its inputs are the original
free-parameter tuple in increasing index order, followed by evidence;
fixed tuple projections supply each child and the evidence Boolean test.
Theorems~\ref{thm:p05-original-sigma-evidence} and
\ref{thm:p05-pa-pr-graphs}, with the record verification in
\cref{annex:p5.original-sigma-evidence}, give, with universal closure,
\[
\begin{aligned}
 \PA&\derives E_\varphi(\vec x,e)\to\exists p\exists q\,\\
 &\quad(\Gamma_{\widehat d_\varphi^+}(\vec x,e,p)\land
   \Gamma_{s_\varphi^+}(\vec x,q)\land
   \mathsf{Prf}^{19}_{\PA}(p,q)).
\end{aligned}
\tag{9}
\]
Here \(s_\varphi^+\) is the same finite numeral-substitution description
defined before (4), now for the fixed original \(\varphi\). To spell
out the bridge to (9), the signed theorem's premise
\(\mathcal E_\varphi^+\) is literally
\(\mathcal E_{\mathcal N^+\varphi}=E_\varphi\). Its safe-tree
conclusion has the original root \(\varphi^\nu\). The bounded-original
clause uses the two builders in (4) and the internally proved signed
Boolean identities (S1)--(S3) in
\cref{annex:p5.original-sigma-evidence}. The other signed clauses build
the original connective or quantifier: in particular, implication evidence
selects a proof of the negative antecedent or positive consequent, and the
output introduces the original implication. Marked collection uses the
actual arbitrary-matrix constructor of
\cref{thm:p05-bounded-case-elimination}, with the original guard and
the complete list of child proofs. Its empty strict and one-entry weak
zero cases are included.

The original numeral-instance root is a sentence. Packing the safe tree
therefore satisfies every factor of the selected \(d_{19}\) checker;
the outer successful evidence guard returns exactly that packed record.
Introduce the unique \(\Gamma_{\widehat d_\varphi^+}\) and
\(\Gamma_{s_\varphi^+}\) outputs to obtain (9). This is the
record-level consequence proved at the end of
\cref{annex:p5.original-sigma-evidence}; the normal form determines the
evidence, while the signed root calculation supplies the original code.

For the closed \(\sigma\), its parameter tuple is empty, so
\(s_\sigma^+\) is the constant \(\ulcorner\sigma\urcorner\) and
\(\widehat d_\sigma^+\) is unary. Eliminating the constant graph
output in (9) gives exactly
\[
 \PA\derives\forall e\,\bigl(E_\sigma(e)\to
  \exists p\,(\Gamma_{\widehat d_\sigma^+}(e,p)\land
   \mathsf{Prf}^{19}_{\PA}(p,\overline{\ulcorner\sigma\urcorner}))\bigr).
\]
Assume \(\sigma\) inside \(\PA\). By (7) choose its evidence
\(e\); the displayed implication gives a proof value \(p\).
Conjunction elimination gives its selected proof assertion. Existential
introduction in \(p\), elimination of the evidence witness, and discharge
of \(\sigma\) give
\[
 \PA\derives
 \sigma\to\exists p\,\mathsf{Prf}^{19}_{\PA}
   (p,\overline{\ulcorner\sigma\urcorner}),
\]
which is the formalized clause
\(\PA\derives\sigma\to
\mathsf{Prv}^{19}_{\PA}(\overline{\ulcorner\sigma\urcorner})\)
by \cref{def:p05-selected-direct-pa-graph}.

Finally suppose externally that \(\mathcal N\entails\sigma\). The fixed
logical equivalence gives \(\mathcal N\entails\mathcal N^+\sigma\).
Structural recursion on this positive formation chooses an ordinary finite
evidence number \(e\): use zero at a bounded leaf, pair at conjunction,
tag a true disjunct, and pair each existential witness with its matrix
evidence, retaining a marked witness's bound. At a bounded universal,
there are only finitely many indices and Lemma~\ref{lem:goedel-beta}
codes their separately chosen evidence objects. At a strict zero bound
there are none; a weak zero bound requires the single object at zero.
Thus \(\mathcal N\entails E_\sigma(\bar e)\).  The positive half of (3) gives
\(\mathsf Q\derives E_\sigma(\bar e)\), and (6) gives
\(\mathsf Q\derives\sigma\).  This proves the external clause and completes
both claims in the body theorem without a prefix-normal-form assumption.
\end{annexitem}
% END HISTORICAL SLICE 16-34
